\newcommand{\fontstxs}{\fontsize{15pt}{17pt} \selectfont \songti{}}
\newcommand{\fontstsh}{\fontsize{14pt}{16pt} \selectfont \songti{}}
\newcommand{\fonthtxe}{\fontsize{18pt}{20pt} \selectfont \heiti{}}

\newcommand{\makecover}{
  \thispagestyle{empty}
  \begin{center}
    \vspace*{1cm}
    \noindent

    %% 这里两个minipage的作用是让Logo与标题对齐

    \begin{minipage}[c]{2.17cm}
      \includegraphics[width=\linewidth]{assets/MUC.jpg}
    \end{minipage}
    \hspace{0.5cm}
    \begin{minipage}[c]{0.8\textwidth}
      {\fontsize{18pt}{20pt}\selectfont \heiti{}
      中央民族大学大学生创新训练计划结题论文\\[0.3cm]}
      {\fontsize{12pt}{14pt}\fontfamily{phv} \selectfont \bfseries
      URTP Final Thesis of Minzu University of China}

      % 中文黑体小二，英文Arial小四加粗
    \end{minipage}

    %% 下方是正式的内容

    \vspace{4cm}
    \makedoctitle{}
    \vspace{0.5cm}
    \makedocsubtitle{}
    \vspace{6cm}

    \begin{tabular}{ll}
      {\fontstxs{} 计划编号：} & \underline{\fontstxs{} \makebox[6cm][c]{\varmucprojectid}}\\
      {\fontstxs{} 项目类型：} & \underline{\fontstxs{} \makebox[6cm][c]{\varmucprojecttype}}\\
      {\fontstxs{} 指导教师：} & \underline{\fontstxs{} \makebox[6cm][c]{\varmucprojectteacher}}\\
      {\fontstxs{} 项目成员：} & \underline{\fontstxs{} \makebox[6cm][c]{\varmucprojectmembers}}\\
      {\fontstxs{} 依托单位：} & \underline{\fontstxs{} \makebox[6cm][c]{\varmucprojectunit}}\\
    \end{tabular}

    \vspace{3cm}

    {\fontstsh{} \varmucprojectdate{}}
  \end{center}
  \clearpage
}